\section{Instance Semantics and Guarantees}
\label{sec:semantics}
The semantic object is an execution domain together with selected observations.
We first define its meaning, then state the contracts that connect symbolic
execution and proof checking to an instance verdict.

\subsection{Instances and observations}
A concrete configuration $\sigma$ contains the control location, C environment,
heap, and call stack. A logical environment $\eta$ assigns values to typed
logical variables. We write $z=(\sigma,\eta)$; $\eta$ retains the root's entry
values throughout execution. The interpretation $\Gamma$ supplies the logical
definitions, memory predicates, and target layout. Assertions are predicates
on $z$, with separating conjunction interpreted by disjoint owned resources
as in separation logic~\cite{reynolds2002separation}.

\begin{definition}[Specification instance]
An instance is $I=\langle c,R,E,\theta,\mathcal B,\Gamma,b\rangle$, where $c$ is
the C body, $R$ and $E$ are its entry and return assertions, $\theta$ is a typed
partial binding, $\mathcal B$ selects observation boundaries, and $b$ bounds
execution and proof-search resources. Its input domain is
\begin{equation}
\dom=\{(\sigma,\eta)\mid\eta\models\theta\land
                         (\sigma,\eta)\models_\Gamma R\}.
\end{equation}
\end{definition}
A binding constrains selected values without choosing all unbound values.
For example, $x=13$ leaves the owned cell's address symbolic. Entry existence
asks whether $\dom\ne\varnothing$. Subsequent execution checking quantifies
over the entire domain. An existence witness therefore does not replace the
domain by one concrete heap. Stronger bindings induce smaller domains, a
property proved in \texttt{InstanceChecking.v}.

An observation resolves the selected assertion against the current C
configuration and its immutable logical entry bindings. Let
$\mathcal B_\Gamma(z)$ denote the finite list of resolved observations at $z$.
For a finite execution trace $\tau=z_0\cdots z_m$, define
\[
\mathsf{Observe}_{\mathcal B}(\tau)
 = \mathop{\bigoplus}_{i=0}^{m}[(z_i,A)\mid A\in\mathcal B_\Gamma(z_i)],
\]
where $\bigoplus$ concatenates lists, retaining order and multiplicity. A checked
trace satisfies every recorded $(z_i,A)$, has safe program transitions, and
satisfies $E$ at its terminal state. Calls introduce fresh logical bindings for
the callee, keeping the caller's bindings available for return observations.

\subsection{Symbolic execution and boundary rules}
The symbolic configuration is $\langle\ell,P,K,U,\Omega,b\rangle$: a C control
location, assertion describing possible states, call/continuation context,
pending observation list, accumulated obligations, and remaining budget.
Annotations are scheduled alongside C locations; consuming an observation
advances $U$ while retaining the C control location. A concrete step relation
$\to_c$ supplies the program semantics. A symbolic step producing successors
$P_1,\ldots,P_k$ must establish progress, safety, and coverage:
\begin{equation}
\label{eq:coverage}
\begin{split}
 &\forall z\in\llbracket P\rrbracket.\;
    \mathsf{Safe}(z)\land\exists z'.\ z\to_c z',\\
 &\forall z,z'.\;z\in\llbracket P\rrbracket\land z\to_c z'
    \Rightarrow z'\in\bigcup_j\llbracket P_j\rrbracket.
\end{split}
\end{equation}
Here and below interpretations use $\Gamma$. Coverage permits an overapproximation for acceptance. The negative-result
interface additionally requires that its admitted path states are realizable. An unresolved branch condition retains all successors whose
infeasibility has not been established.

At an observation, the administrative rule is
\begin{equation}
\label{eq:observe}
\frac{Q=\mathsf{Resolve}_{\Gamma,K}(A)*F_{P,A}}
 {\langle P,A::U,\Omega\rangle
  \xrightarrow{\mathrm{observe}}
  \langle P,U,\Omega\cup\{P\vdash Q\}\rangle}.
\end{equation}
The omitted context $(\ell,K,b)$ and represented program state $P$ remain unchanged. The exact local-storage frame $F_{P,A}$ accounts for live
scalar stack cells outside the observation's footprint
(Section~\ref{sec:boundaries}). Crucially, $P\vdash Q$ is recorded as an
obligation, rather than required to enable the continuation.

The other administrative rules have explicit interfaces. Call entry copies
evaluated arguments to fresh parameter slots, saves the caller context, and
enters the actual body. Return retains the actual return value and shared
writes, retires callee-local storage, restores the caller context, and records
the selected callee postcondition on a checking copy. Loops repeatedly apply
the condition and body transitions, recording observations at visited invariant
locations. An exhausted execution budget creates an unfinished leaf.

\subsection{Observation noninterference}
\label{sec:noninterference}
The reference observation semantics records obligations independently of their
truth. Its state already includes the program control location, so a boundary
schedule can distinguish entries, returns, and individual loop visits.
\texttt{ObservationTrace.v} defines ordinary finite executions and monitored
executions with the same state trace and an additional observation list.

\begin{theorem}[Erasure and lifting]
\label{thm:erasure}
For any resolved finite boundary schedule $\mathcal B$ and finite trace $\tau$
from $z$ to $z'$, there exists a monitored execution with state trace $\tau$
if and only if $\tau$ is a program execution from $z$ to $z'$. Its observation
list is exactly $\mathsf{Observe}_{\mathcal B}(\tau)$.
\end{theorem}
\begin{proof}
Induct on the finite execution derivation. A zero-step trace records the
observations of its sole state. A step preserves the concrete transition,
prepends that state's observations, and applies the induction hypothesis to
the remaining trace. Conversely, erasing the observation list from a monitored
derivation preserves every program transition. No constructor requires a
recorded assertion to hold. List induction establishes the observation inventory.
\end{proof}
Thus changing an assertion's truth cannot prune or create a program trace in
the reference model. If all recorded obligations are discharged, each selected
assertion holds at each corresponding visited state. Both directions and the observation-coverage property are mechanized.

\subsection{Acceptance and negative evidence}
\label{sec:verdicts}
A finite symbolic execution tree has program-step, normalization, observation,
terminal, and unfinished nodes. Terminal nodes require safety, the return
assertion, and the absence of further program transitions. Program-step nodes
satisfy Equation~\eqref{eq:coverage}; observation nodes follow
Equation~\eqref{eq:observe}. Normalization preserves the represented joint
observations (Section~\ref{sec:normalization}). An unfinished node has no
acceptance certificate.

\begin{samepage}
\begin{theorem}[Nonvacuous instance acceptance]
\label{thm:acceptance}
If the root represents $\dom$, the domain is nonempty, and every obligation
of a finite covering execution tree is valid, then every initial state in
$\dom$ has safe, terminating executions satisfying $E$. Under the boundary
coverage interface, every selected observation on those executions also holds.
\end{theorem}
\end{samepage}
\begin{proof}
Induction on the tree establishes safe termination. A step covers every
concrete successor and applies the induction hypothesis to each covering
subtree. Normalization transfers membership through its semantic equivalence;
observation preserves the continuation state. Terminal obligations establish
the postcondition. Entry existence supplies nonvacuity independently. For
observations, Theorem~\ref{thm:erasure}'s inventory and checked boundary
entailments establish satisfaction at each selected visit.
\end{proof}
The execution and nonvacuity result is mechanized in
\texttt{InstanceChecking.v}; observation coverage is proved separately in
\texttt{ObservationTrace.v}. The implementation obligations corresponding to these interfaces are collected
in Section~\ref{sec:evidence}.

A confirmed violation has an existential meaning:
\begin{equation}
\label{eq:violation}
\exists z_0\in\dom, z.\;z_0\to_c^*z\land\mathsf{Bad}_{E,\mathcal B}(z).
\end{equation}
Here badness is an unsafe state, a violated selected observation, or a terminal
state violating $E$. The test executor preserves the entry constraints, branch
conditions and actual state updates. For a path VC $P\vdash A$, SMT/ATP may
discharge routine positive obligations or supply a candidate satisfying
$P\land\neg A$. The agent may attempt either a proof of the original VC or a
proof of its negation; both are checked against the frozen obligation. A
checked positive proof establishes \valid{}. A checked negative proof with an
admitted reachable path witness establishes \invalid{}. If neither direction
is established, the VC remains \unknown{}. Infeasible paths cannot cause a
failed test merely because an implication proof was unfinished.

\begin{table}[t]
\centering\small
\caption{Entry results, VC states, and whole-test verdicts.}
\label{tab:verdicts}
\begin{tabular}{@{}p{.27\columnwidth}p{.68\columnwidth}@{}}
\toprule
Layer & Result and meaning\\
\midrule
Entry & SAT: the instance domain is nonempty. UNSAT: it is empty. UNKNOWN: satisfiability is unresolved.\\
VC & \valid{}: established. \invalid{}: an admitted path check establishes a violation. \unknown{}: unresolved.\\
Test: \pass & Entry SAT, complete execution, and all selected obligations VALID.\\
Test: \fail & Entry UNSAT, or a selected path VC INVALID; other obligations may remain unresolved.\\
Test: \unknown & Neither complete acceptance nor a confirmed failed check.\\
\bottomrule
\end{tabular}
\end{table}
Table~\ref{tab:verdicts} separates logical results from execution bookkeeping.
The report exposes \texttt{precondition\_result}, per-obligation
\texttt{vc\_results}, and \texttt{test\_result}. An UNSAT entry contributes
a false result to the conjunction and stops execution. Infrastructure errors
are reported separately. Legacy run-status fields remain available for
compatibility. Unsupported logical encodings retain their original obligations
for proof completion. Execution correspondence is part of the executor's
trusted interface: an overapproximated proof-search state alone cannot license
an INVALID path result.

The generated initial VC uses a negative-entailment
encoding; its equivalence to existence is proved under an explicit classical
logic premise in \texttt{PreconditionOracle.v}. Execution-only and
Require-only checks retain different scopes even when both are reported as
\pass{} by the command-line interface.
